\begin{scanpage}{0209}{191}
\noindent the semi-infinite conductor $(-\infty,z_2)$ as a plane wave. Then the current induced in this section of the conductor will be determined by the equation
\[J(z)=S\bigl[e^{iwz}-e^{iwz_2}\psi_+(z_2-z)e^{ik(z_2-z)}\bigr],\tag{30.26}\]
where
\[S=\frac{i\omega E'_{0z}}{2k^2\sin^2\vartheta\ln\dfrac{2i}{\gamma ka\sin\vartheta}},\quad w=-k\cos\vartheta.\tag{30.27}\]
We will select the quantity $z_0$ in such a way that, at a distance $z_2-z_0$ from the conductor end, the reflected current wave would be practically equal to zero $(\psi_+(z_2-z_0)\approx0)$. Substituting the function (30.26) into the right-hand member of Equality (30.25) and taking for the quantity $E_z^e$ the external field (30.16), we obtain
\[
\begin{aligned}
E_{2\vartheta}=H_{2\varphi}={}&\frac{1}{2\sin\vartheta\ln\dfrac{2i}{\gamma ka\sin\vartheta}}\times\\
&\times\frac{e^{ikR}}R\int_{-\infty}^{z_2}E_z^e\bigl[e^{iwz}-\psi_+(z_2-z)e^{iwz_2+ik(z_2-z)}\bigr]dz\\
={}&\frac{1}{2\sin\vartheta\ln\dfrac{2i}{\gamma ka\sin\vartheta}}\frac{e^{ikR}}R\Biggl\{\mathcal E_1\bigl[e^{iwz_1}-\psi_+(L)e^{iwz_2+ikL}\bigr]\\
&+\hat E_{0z}\frac{e^{i(w+w_0)z_1}}{i(w+w_0)}
-\hat E_{0z}e^{i(k+w)z_2}\int_{-\infty}^{z_1}e^{-i(k-w_0)z}\psi_+(z_2-z)\,dz\Biggr\}.
\end{aligned}\tag{30.28}
\]
An important feature of this relationship is that the integration is performed here not along the entire conductor $(-\infty,z_2)$, but only along part of it $(-\infty,z_1)$, where the function $\psi_+(z_2-z)$ describes the current with good precision. The integrals here are calculated the same as in Equation (30.09). As a result, the field radiated by the semi-infinite conductor $(-\infty,z_2)$ will equal
\[E_{2\vartheta}=H_{2\varphi}=\frac{1}{2\sin\vartheta\ln\dfrac{2i}{\gamma ka\sin\vartheta}}\frac{e^{ikR}}R\times\]
\begin{flushright}(Equation continued on next page)\end{flushright}
\end{scanpage}
